% -------------------------------------------------------------------------
% ÜBERHOLT (31.08.2026). Maßgeblich ist die eingereichte main.tex. Zahlen,
% Zeilennummern und Angaben zum Overleaf-Stand in dieser Datei stammen aus
% der Zeit VOR der Crosswalk-Korrektur vom 29.08.2026 (35 statt 36 Stadt-
% teile, vollständiger Neulauf) und wurden bewusst NICHT nachgezogen: die
% Datei ist als datierter Entwurf Teil der Arbeitsdokumentation.
% -------------------------------------------------------------------------
% =========================================================================
% KAPITEL 5.2 -- ZUSAMMENFUEHRUNG DER DATENQUELLEN
% Fassung 1, 24.08.2026. STAND IN OVERLEAF -- eingetippt und validiert.
%
% UEBERSCHRIFT: Die Zahl ist gestrichen. Frueher hiess der Abschnitt
% "Zusammenfuehrung der VIER Datenquellen" -- Kapitel 4.1 fuehrt aber sieben
% Quellen. Beide Zaehlungen sind fuer sich richtig, nebeneinander sehen sie
% nach einem Fehler aus.
%
% GLEITOBJEKT: Quelltext 1 (lst:acs-versatz), liegt in
% docs/kapitel5-2_gleitobjekte.tex. Direkt nach Absatz 2 einfuegen.
%
% KEINE FUSSNOTE IN DIESEM ABSCHNITT -- und das ist Absicht. Alle sieben
% Quellen sind in 4.1 belegt, der Publikationsversatz in 4.2 mit
% CensusBureau2020. Ein zweiter \footcite auf dieselbe Seite waere die
% Dopplung zwischen Grundlagen- und Empirieteil, die R8 abstraft.
%
% ZWEI SETZUNGEN, DIE ZU KAPITEL 4 PASSEN MUESSEN
%   - Der Kriminalitaetsschnitt ist auf Januar 2018 datiert, wie in 4.1
%     ("monatlich verdichtete Fassung ab Januar 2018", "Altquelle bis
%     Dezember 2017"). Das Datum des eigentlichen Systemwechsels (Mai 2018)
%     steht hier NICHT -- es traegt in 5.3 das Argument fuer den relativen
%     Index und gehoert an eine Stelle, nicht an zwei.
%   - 154.544 ist die Zahl der ZUGEORDNETEN Parzellen, nicht die Rohzahl.
%     Die Rohzahl 155.395 steht in Kapitel 4, weil der Join dort noch nicht
%     stattgefunden hat.
%
% amsmath wird fuer \text{} in der Formel gebraucht. Ohne amsmath tut es
% \mathrm{} genauso.
% =========================================================================


% ---------------------------- AB HIER KOPIEREN ---------------------------

\subsection{Zusammenführung der Datenquellen}
\label{sec:integration}

Die sieben Quellen aus Abschnitt~\ref{sec:datengrundlage} liegen in drei
Raumbezügen und zwei Zeitrastern vor: Einsätze mit Stadtteilangabe und
Zeitstempel, sozioökonomische Merkmale je Census Tract und Jahrgang, bauliche
Merkmale je Parzelle ohne Zeitbezug. Sie werden auf der Ebene des einzelnen
Einsatzes zusammengeführt, nicht auf der späteren Analyseeinheit. Der Grund ist
der Zeitbezug: Jeder Einsatz trägt seinen eigenen Zeitpunkt, und nur auf dieser
Ebene lässt sich jedem Einsatz derjenige Stand der zeitabhängigen Merkmale
zuordnen, der damals verfügbar war. Verdichtet wird erst danach
(Abschnitt~\ref{sec:merkmale}); die Reihenfolge folgt damit der Verarbeitung
und nicht der Aufgabenliste des Vorgehensmodells, das die Konstruktion vor die
Zusammenführung stellt. Das Ergebnis ist eine Tabelle mit 719.989 Zeilen und
50 Spalten, eine je Einsatz.

Die sozioökonomischen Merkmale werden in zwei Stufen angefügt. Räumlich ordnet
ein Crosswalk jeden Census Tract einem Stadtteil zu; dabei werden Zählgrößen
summiert, Medianwerte dagegen bevölkerungsgewichtet gemittelt, weil sich ein
Median nicht addieren lässt. Zeitlich gilt die Bedingung
$\text{ACS-Jahrgang} \le \text{Einsatzjahr} - 1$: Einem Einsatz wird der
jüngste Jahrgang zugeordnet, der zu seinem Zeitpunkt bereits veröffentlicht
war. Ohne diese Bedingung erhielte ein Einsatz des Jahres 2023 den Jahrgang
2023, der erst gegen Ende 2024 erschien (Abschnitt~\ref{sec:eda}).
Quelltext~\ref{lst:acs-versatz} zeigt die Auswahl. Vollständig ist die
räumliche Zuordnung nicht: Der Crosswalk beruht auf den Tract-Grenzen des
Zensus 2020, während die älteren Jahrgänge abweichende Grenzen verwenden.
Zugeordnet werden 63,1\,\% der Tracts des Jahrgangs 2009, 79,7\,\% der
Jahrgänge 2014 und 2019 und 99,2\,\% der Jahrgänge 2021 und 2023.

Die Kriminalitätsbelastung stammt aus zwei Quellen, deren Schnitt im Januar
2018 liegt. Die Altquelle führt keine Stadtteilspalte, dafür Koordinaten; ihre
Meldungen werden deshalb über einen räumlichen Verbund demjenigen
Stadtteilpolygon zugeordnet, in dem sie liegen. Verwendet wird dabei dieselbe
Geometrie wie für die baulichen Merkmale, weil sich Kriminalitäts- und
Baumerkmale desselben Stadtteils sonst auf unterschiedliche Flächen bezögen.
Der Schnitt wird am Beginn der Nachfolgequelle gezogen und nicht am Ende der
Altquelle, damit kein Monat aus zwei Erfassungssystemen gespeist wird. Fällt
die Zuordnungsquote des räumlichen Verbunds unter 90\,\%, bricht die
Aufbereitung ab.

Die baulichen Merkmale liegen je Parzelle vor. Jede Parzelle wird über ihren
Flächenmittelpunkt dem Stadtteil zugeordnet, in dem dieser liegt; das gelingt
für 154.544 Parzellen und damit für 99,5\,\% des Verzeichnisses. Der Bestand
ist ein Abzug des Jahres 2020 und der einzige verfügbare Jahrgang. Die
baulichen Merkmale sind deshalb je Stadtteil über den gesamten Analysezeitraum
konstant und tragen keine zeitliche Variation bei (Abschnitt~\ref{sec:eda}).

Alle drei Zusammenführungen sind nach demselben Muster abgesichert. Vor jedem
Verbund wird geprüft, ob der Schlüssel auf der angefügten Seite eindeutig ist;
nach jedem Verbund wird die Zeilenzahl gegen den Erwartungswert gehalten, und
bei Abweichung bricht die Aufbereitung ab. Diese Prüfung ist nicht
Sorgfaltsroutine, sondern verhindert einen Fehler, der sich sonst nicht
bemerkbar macht: Ist der Schlüssel mehrdeutig, entsteht ein kartesisches
Produkt, die Tabelle enthält mehr Zeilen als Einsätze, und jeder betroffene
Einsatz geht mehrfach in die spätere Aggregation ein.

% ---------------------------- BIS HIER KOPIEREN --------------------------
